\subsubsection{Fidelity of a Quantum Gate}\label{sec:fidelity-of-a-quantum-gate}

We have defined the types of errors that can occur in quantum gates. But \textbf{\emph{how can we measure how well a quantum gate performs?}} The \textbf{fidelity} of a quantum gate is the answer to this question. It is a measure of how close the actual operation performed by the gate is to the ideal operation.

\highspace
\begin{flushleft}
    \textcolor{Green3}{\faIcon{question-circle} \textbf{Motivation}}
\end{flushleft}
In an ideal quantum computer, every gate would be implemented perfectly. If we want to apply a gate $U_I$ (where $I$ stands for ``\emph{ideal}''), then the \textbf{output should be exactly}:
\begin{equation}\label{eq:ideal-output}
    y_I = U_I x
\end{equation}
Where $x$ is the input qubit and the $y_I$ is the \emph{ideal} output. Unfortunately, real hardware is noisy (as we have seen in the previous sections). The gate physically implemented on the processor is slightly different from the ideal one: $U_I \neq U_P$ (where $P$ stands for ``\emph{physical}''). Therefore, the \textbf{actual output} is:
\begin{equation}\label{eq:actual-output}
    y_P = U_P x
\end{equation}
This output may differ from the ideal one due to imperfect control electronics, calibration errors, decoherence, and other sources of noise like environmental interactions. Given this situation, we ask ourselves: \hl{how close is $y_P$ to $y_I$?} How closely does the physical output $y_P$ produced by the real gate $U_P$ resemble the ideal output $y_I$ we would get from a perfect gate $U_I$? We introduce the \textbf{fidelity} to answer this question.

\highspace
\begin{definitionbox}[: Fidelity]
    The \definition{Fidelity} of a quantum gate is a measure of \textbf{how closely the actual operation performed by the gate matches the ideal operation}. It is defined as:
    \begin{equation}\label{eq:fidelity}
        F = \left| \braket{y_I}{y_P} \right|^2
    \end{equation}
    Where $\ket{y_I}$ is the ideal output state and $\ket{y_P}$ is the actual output state produced by the physical gate. We use the inner product $\braket{y_I}{y_P}$ to quantify the \textbf{overlap between the two states}, and the \textbf{square} of its magnitude gives us a \textbf{value between 0 and 1} that represents the fidelity.
    The fidelity ranges from 0 to 1, where:
    \begin{itemize}
        \item $F = 1$ means the \textbf{actual} output \textbf{perfectly matches the ideal} output (i.e., the \hl{gate is perfect}).
        \item $F = 0$ means the \textbf{actual} output is \textbf{completely orthogonal to the ideal} output (i.e., the \hl{gate is completely wrong}).
        \item $0 < F < 1$ indicates that the \textbf{actual} output is \textbf{partially close} to the ideal output, with higher values indicating better performance (i.e., the \hl{gate is somewhat good}).
    \end{itemize}

    We can also write the fidelity more explicitly in terms of the gates. Using the \autoref{eq:ideal-output} and \autoref{eq:actual-output}:
    \begin{equation*}
        y_I = U_I x \quad \text{and} \quad y_P = U_P x
    \end{equation*}
    Writing the conjugate transpose of $y_I$ (bra):
    \begin{equation*}
        \ket{y_I} = U_I \ket{x} \quad \Rightarrow \quad \bra{y_I} = \bra{x} U_I^\dagger
    \end{equation*}
    Substituting this into the fidelity definition:
    \begin{equation}\label{eq:fidelity-explicit}
        F = \left| \braket{y_I}{y_P} \right|^2 = \left| \bra{x} U_I^\dagger U_P \ket{x} \right|^2
    \end{equation}
\end{definitionbox}

\highspace
\begin{flushleft}
    \textcolor{Red2}{\faIcon{exclamation-triangle} \textbf{Why isn't one fidelity value enough?}}
\end{flushleft}
Suppose we build a terrible \textbf{physical gate}:
\begin{equation*}
    U_P = I = \begin{bmatrix}
        1 & 0 \\
        0 & 1
    \end{bmatrix}
\end{equation*}
This gate does nothing to the input state. Now suppose the \textbf{ideal gate} should be:
\begin{equation*}
    U_I = X = \begin{bmatrix}
        0 & 1 \\
        1 & 0
    \end{bmatrix}
\end{equation*}
Let's test the fidelity using only \ket{+}:
\begin{equation*}
    U_I\ket{+} = X\ket{+} = \ket{+}
\end{equation*}
Because \ket{+} is an eigenstate of $X$. Meanwhile:
\begin{equation*}
    U_P\ket{+} = I\ket{+} = \ket{+}
\end{equation*}
Calculating the fidelity:
\begin{equation*}
    F = \left| \bra{+} X^\dagger I \ket{+} \right|^2 = \left| \bra{+} X I \ket{+} \right|^2 = \left| \bra{+} X \ket{+} \right|^2 = \left| \braket{+}{+} \right|^2 = 1
\end{equation*}
Perfect fidelity! But we know that $U_P$ is a terrible gate. In fact, if we test the fidelity with \ket{0}:
\begin{equation*}
    U_I\ket{0} = X\ket{0} = \ket{1}
\end{equation*}
While:
\begin{equation*}
    U_P\ket{0} = I\ket{0} = \ket{0}
\end{equation*}
Calculating the fidelity:
\begin{equation*}
    F = \left| \bra{0} X^\dagger I \ket{0} \right|^2 = \left| \bra{0} X I \ket{0} \right|^2 = \left| \bra{0} X \ket{0} \right|^2 = \left| \braket{0}{1} \right|^2 = 0
\end{equation*}
\textbf{Terrible fidelity!} The \hl{same gate appears perfect and terrible depending on which input state we choose.} So asking ``\emph{what is the fidelity of this gate?}'' is incomplete, instead we should ask ``\emph{for which input state?}''.

\highspace
\begin{flushleft}
    \textcolor{Green3}{\faIcon{check-circle} \textbf{The Solution: Average Fidelity}}
\end{flushleft}
Instead of testing a single input state, we conceptually test \textbf{all possible input states}. Remember:
\begin{equation*}
    \ket{\psi} = \alpha \ket{0} + \beta \ket{1}
\end{equation*}
Can be any point on the Bloch sphere. So rather than computing the fidelity for a single state (\autoref{eq:fidelity-explicit}):
\begin{equation*}
    F(\ket{\psi}) = \left| \bra{\psi} U_P^{\dagger} U_I \ket{\psi} \right|^2
\end{equation*}
For one specific \ket{\psi}, we \textbf{average over all possible input states}:
\begin{equation}
    F_{\text{avg}} = \int \left| \bra{\psi} U_P^{\dagger} U_I \ket{\psi} \right|^2 \, \mathrm{d}\psi
\end{equation}
This is the \definition{Average Fidelity} of the quantum gate. The notation $\mathrm{d}\psi$ indicates a \textbf{uniform average over all possible quantum states}. For a single qubit, every pure state can be represented as a point on the Bloch sphere. Therefore, the average fidelity \hl{measures the performance of the gate over the entire Bloch sphere rather than for a specific input state.}

\highspace
In other words, we conceptually:
\begin{enumerate}
    \item Choose a quantum state $\ket{\psi}$;
    \item Apply both the ideal gate $U_I$ and the physical gate $U_P$;
    \item Compute the fidelity;
    \item Repeat the process for all possible states on the Bloch sphere;
    \item Average the results.
\end{enumerate}
The resulting value provides a single number that characterizes the overall quality of the gate, independently of the particular input state chosen.

\highspace
\textcolor{Green3}{\faIcon{question-circle} \textbf{Can we really test all possible states?}} \hl{In practice}, this is \textbf{impossible because the Bloch sphere contains infinitely many states}. Fortunately, it can be shown (proof outside the scope of this course) that the average fidelity can be computed directly from the matrices describing the ideal and physical gates:
\begin{equation}\label{eq:average-fidelity}
    F_{\text{avg}}
    =
    \frac{
        \left|
        \mathrm{Tr}
        \left(
            U_P^{\dagger}U_I
        \right)
        \right|^2
        + d
    }
    {
        d(d+1)
    }
\end{equation}
Where Tr is the \textbf{trace of a matrix}, a sum of the diagonal entries:
\begin{equation*}
    \mathrm{Tr}(A) = \sum_i A_{ii}
\end{equation*}
And:
\begin{equation*}
    d = 2^n
\end{equation*}
Is the \textbf{dimension of the Hilbert space} and $n$ is the \textbf{number of qubits} involved in the gate.

\highspace
\textcolor{Green3}{\faIcon[regular]{lightbulb} \textbf{Meaning of $d$.}} For a single-qubit gate:
\begin{equation*}
    d = 2^1 = 2
\end{equation*}
Because a qubit lives in a two-dimensional Hilbert space. For a two-qubit gate:
\begin{equation*}
    d = 2^2 = 4
\end{equation*}
Because the computational basis contains four states:
\begin{equation*}
    \ket{00},\;
    \ket{01},\;
    \ket{10},\;
    \ket{11}
\end{equation*}
More generally, an $n$-qubit system has
\begin{equation*}
    2^n
\end{equation*}
Basis states, hence the dimension of its Hilbert space is
\begin{equation*}
    d = 2^n
\end{equation*}

\highspace
\textcolor{Green3}{\faIcon{check-circle} \textbf{Interpretation.}} \hl{The average fidelity replaces an impossible average over infinitely many quantum states with a simple calculation involving only the matrices of the ideal and physical gates.} This allows hardware vendors and researchers to assign a single fidelity value to a gate and compare different quantum processors.